% ============================================================================
% ERWARTUNGSHORIZONT - Stundenleistung: Bewegungen
% ============================================================================
% Lösungen in ROT, mit Punkteverteilung
% ============================================================================

\documentclass[11pt,a4paper]{article}

\usepackage[utf8]{inputenc}
\usepackage[T1]{fontenc}
\usepackage[ngerman]{babel}
\usepackage{geometry}
\usepackage{helvet}
\usepackage{tikz}
\usepackage{amsmath}
\usepackage{amssymb}
\usepackage{enumitem}
\usepackage{tabularx}
\usepackage{array}
\usepackage{fancyhdr}
\usepackage{lastpage}
\usepackage{xcolor}

\renewcommand{\familydefault}{\sfdefault}

\geometry{left=2cm, right=2cm, top=2.8cm, bottom=2cm}

% Farben
\definecolor{physik}{HTML}{2c5aa0}
\definecolor{lightblue}{HTML}{e3f2fd}
\definecolor{lightgray}{HTML}{f5f5f5}
\definecolor{loesung}{HTML}{c62828}

% Kopfzeile
\pagestyle{fancy}
\fancyhf{}
\fancyhead[L]{\small\textcolor{loesung}{\textbf{ERWARTUNGSHORIZONT}} | Physik Klasse 6 | Bewegungen}
\fancyhead[R]{\small Stundenleistung Stunde 4}
\fancyfoot[C]{\small Seite \thepage{} von \pageref{LastPage}}
\renewcommand{\headrulewidth}{0.4pt}

% Lösung-Befehl
\newcommand{\lsg}[1]{\textcolor{loesung}{\textbf{#1}}}

% Lücke mit Lösung
\newcommand{\lueckelsg}[1]{\underline{\lsg{\,#1\,}}}

% Niveaustufen
\newcommand{\stern}[1]{\textbf{(#1)}}

% Punktefeld
\newcommand{\punkte}[1]{\hfill \textit{/#1 P}}

\begin{document}

% ============================================================================
% TITEL
% ============================================================================
\begin{center}
    {\Large\bfseries\textcolor{physik}{Stundenleistung: Bewegungen}}\\[0.3em]
    {\normalsize\textcolor{loesung}{\textbf{ERWARTUNGSHORIZONT}} | Physik Klasse 6}
\end{center}

\vspace{0.5em}
\hrule
\vspace{0.8em}

% ============================================================================
% AUFGABE 1: Bewegungsarten (*) - 8 Punkte
% ============================================================================
\noindent\textbf{Aufgabe 1 \stern{*}: Bewegungsarten zuordnen} \punkte{8}

\vspace{0.5em}
\begin{center}
\renewcommand{\arraystretch}{1.4}
\begin{tabular}{|l|c|c|c||c|c|}
\hline
\textbf{Bewegung} & \multicolumn{3}{c||}{\textbf{Bahnform}} & \multicolumn{2}{c|}{\textbf{Bewegungsart}} \\
\cline{2-6}
& geradlinig & kreisförmig & krummlinig & gleichförmig & ungleichförmig \\
\hline
Rolltreppe (konstante Fahrt) & \lsg{X} & & & \lsg{X} & \\
\hline
Karussell & & \lsg{X} & & \lsg{X} & \\
\hline
Achterbahn & & & \lsg{X} & & \lsg{X} \\
\hline
100-m-Sprinter (Start) & \lsg{X} & & & & \lsg{X} \\
\hline
\end{tabular}
\end{center}

\vspace{0.3em}
\lsg{Bewertung: Je korrektes Kreuz 1 P (max. 8 P)}

\vspace{1em}

% ============================================================================
% AUFGABE 2: Geschwindigkeit berechnen (*) - 6 Punkte
% ============================================================================
\noindent\textbf{Aufgabe 2 \stern{*}: Geschwindigkeit berechnen} \punkte{6}

\vspace{0.5em}
\noindent Ein Radfahrer legt in 10 Sekunden eine Strecke von 50 Metern zurück.

\vspace{0.3em}
\noindent\textbf{a)} Gegeben und Gesucht: \hfill \lsg{2 P}

\vspace{0.3em}
\hspace{1.5em} Gegeben: $s =$ \lueckelsg{50 m} \hspace{2em} $t =$ \lueckelsg{10 s}

\vspace{0.3em}
\hspace{1.5em} Gesucht: \lueckelsg{Geschwindigkeit $v$}

\vspace{0.5em}
\noindent\textbf{b)} Berechnung: \hfill \lsg{4 P}

\vspace{0.3em}
\hspace{1.5em} $v = \dfrac{s}{t} = \dfrac{\lueckelsg{50 m}}{\lueckelsg{10 s}} =$ \lueckelsg{5 m/s}

\vspace{0.3em}
\lsg{Formel: 1 P | Einsetzen: 2 P | Ergebnis mit Einheit: 1 P}

\vspace{1em}

% ============================================================================
% AUFGABE 3: Werte ablesen (*) - 6 Punkte
% ============================================================================
\noindent\textbf{Aufgabe 3 \stern{*}: Werte aus dem Diagramm ablesen} \punkte{6}

\vspace{0.5em}
\begin{center}
\begin{tikzpicture}[scale=0.75]
    % Koordinatensystem
    \draw[->] (0,0) -- (7.5,0) node[right] {$t$ in s};
    \draw[->] (0,0) -- (0,5) node[above] {$s$ in m};

    % Skalierung x
    \foreach \x/\label in {0/0, 1/2, 2/4, 3/6, 4/8, 5/10, 6/12} {
        \draw (\x,-0.1) -- (\x,0.1);
        \node[below] at (\x,-0.2) {\small \label};
    }

    % Skalierung y
    \foreach \y/\label in {0/0, 1/10, 2/20, 3/30, 4/40} {
        \draw (-0.1,\y) -- (0.1,\y);
        \node[left] at (-0.2,\y) {\small \label};
    }

    % Gerade
    \draw[thick, physik] (0,0) -- (6,4);

    % Lösungs-Markierungen
    \draw[dashed, loesung] (3,0) -- (3,2) -- (0,2);
    \fill[loesung] (3,2) circle (2pt);
    \node[loesung, above right] at (3,2) {\scriptsize (6|20)};

    \draw[dashed, loesung] (6,0) -- (6,4) -- (0,4);
    \fill[loesung] (6,4) circle (2pt);
    \node[loesung, above left] at (6,4) {\scriptsize (12|40)};

    \draw[dashed, loesung] (1.5,0) -- (1.5,1) -- (0,1);
    \fill[loesung] (1.5,1) circle (2pt);
    \node[loesung, above] at (1.5,1) {\scriptsize (3|10)};

    \node[physik, right] at (5.5,3.5) {\small Fußgänger};
\end{tikzpicture}
\end{center}

\begin{enumerate}[label=\alph*), leftmargin=2em]
    \item Nach $t = 6$ s: $s =$ \lueckelsg{20 m} \hfill \lsg{2 P}
    \item Für $s = 40$ m: $t =$ \lueckelsg{12 s} \hfill \lsg{2 P}
    \item Nach $t = 3$ s: $s =$ \lueckelsg{10 m} \hfill \lsg{2 P}
\end{enumerate}

\vspace{1em}

% ============================================================================
% AUFGABE 4: Diagramm zeichnen (**) - 8 Punkte
% ============================================================================
\noindent\textbf{Aufgabe 4 \stern{**}: Ein s(t)-Diagramm zeichnen} \punkte{8}

\vspace{0.5em}
\noindent Wertetabelle Jogger:

\begin{center}
\begin{tabular}{|c|c|c|c|c|c|}
\hline
Zeit $t$ in s & 0 & 4 & 8 & 12 & 16 \\
\hline
Weg $s$ in m & 0 & 20 & 40 & 60 & 80 \\
\hline
\end{tabular}
\end{center}

\noindent\textbf{a)} Diagramm: \hfill \lsg{6 P}

\begin{center}
\begin{tikzpicture}[scale=0.65]
    % Koordinatensystem
    \draw[->] (0,0) -- (9,0) node[right] {$t$ in s};
    \draw[->] (0,0) -- (0,5.5) node[above] {$s$ in m};

    % Skalierung x
    \foreach \x/\label in {0/0, 2/4, 4/8, 6/12, 8/16} {
        \draw (\x,-0.1) -- (\x,0.1);
        \node[below] at (\x,-0.3) {\small \label};
    }

    % Skalierung y
    \foreach \y/\label in {0/0, 1/20, 2/40, 3/60, 4/80} {
        \draw (-0.1,\y) -- (0.1,\y);
        \node[left] at (-0.3,\y) {\small \label};
    }

    % Hilfslinien
    \draw[lightgray, very thin] (0,1) -- (8,1);
    \draw[lightgray, very thin] (0,2) -- (8,2);
    \draw[lightgray, very thin] (0,3) -- (8,3);
    \draw[lightgray, very thin] (0,4) -- (8,4);
    \draw[lightgray, very thin] (2,0) -- (2,4);
    \draw[lightgray, very thin] (4,0) -- (4,4);
    \draw[lightgray, very thin] (6,0) -- (6,4);
    \draw[lightgray, very thin] (8,0) -- (8,4);

    % LÖSUNG
    \draw[thick, loesung] (0,0) -- (8,4);
    \fill[loesung] (0,0) circle (4pt);
    \fill[loesung] (2,1) circle (4pt);
    \fill[loesung] (4,2) circle (4pt);
    \fill[loesung] (6,3) circle (4pt);
    \fill[loesung] (8,4) circle (4pt);

    \node[loesung, right] at (7,3.5) {\small Jogger};
\end{tikzpicture}
\end{center}

\lsg{Bewertung: 5 Punkte korrekt (5 P) + Gerade verbunden (1 P)}

\vspace{0.5em}
\noindent\textbf{b)} Geschwindigkeit: \hfill \lsg{2 P}

\vspace{0.3em}
\hspace{1.5em} $v = \dfrac{s}{t} = \dfrac{80\,\text{m}}{16\,\text{s}} =$ \lueckelsg{5 m/s}

\lsg{(Alternativ: $v = 20/4 = 5$ m/s oder andere korrekte Kombination)}

\vspace{1em}

% ============================================================================
% AUFGABE 5: Bewegungen vergleichen (**) - 5 Punkte
% ============================================================================
\noindent\textbf{Aufgabe 5 \stern{**}: Bewegungen vergleichen} \punkte{5}

\vspace{0.5em}
\begin{center}
\begin{tikzpicture}[scale=0.65]
    % Koordinatensystem
    \draw[->] (0,0) -- (7,0) node[right] {$t$ in s};
    \draw[->] (0,0) -- (0,5) node[above] {$s$ in m};

    % Skalierung x
    \foreach \x/\label in {0/0, 1/2, 2/4, 3/6, 4/8, 5/10} {
        \draw (\x,-0.1) -- (\x,0.1);
        \node[below] at (\x,-0.2) {\small \label};
    }

    % Skalierung y
    \foreach \y/\label in {0/0, 1/10, 2/20, 3/30, 4/40} {
        \draw (-0.1,\y) -- (0.1,\y);
        \node[left] at (-0.2,\y) {\small \label};
    }

    % A: steil
    \draw[thick, red] (0,0) -- (2.5,4);
    \node[red, right] at (2,3.5) {\small A};

    % B: flach
    \draw[thick, physik] (0,0) -- (5,4);
    \node[physik, right] at (4.5,3.5) {\small B};

    % Lösungs-Markierung
    \draw[dashed, loesung] (2.5,0) -- (2.5,4);
    \node[loesung] at (2.5,-0.5) {\scriptsize 5 s};
\end{tikzpicture}
\end{center}

\begin{enumerate}[label=\alph*), leftmargin=2em]
    \item Schneller: \lueckelsg{A} \hfill \lsg{1 P}

    \item Begründung: \hfill \lsg{2 P}

          \lsg{``Die Gerade von A ist steiler als die von B.''}

          \lsg{oder: ``A legt in gleicher Zeit mehr Weg zurück.''}

          \lsg{oder: ``A erreicht 40 m schneller als B.''}

    \item $t =$ \lueckelsg{5 s} \hfill \lsg{2 P}
\end{enumerate}

\vspace{1em}

% ============================================================================
% AUFGABE 6: Transfer (***) - 6 Punkte
% ============================================================================
\noindent\textbf{Aufgabe 6 \stern{***}: Bewegungsgeschichte interpretieren} \punkte{6}

\vspace{0.5em}
\begin{center}
\begin{tikzpicture}[scale=0.55]
    % Koordinatensystem
    \draw[->] (0,0) -- (11,0) node[right] {$t$ in min};
    \draw[->] (0,0) -- (0,6) node[above] {$s$ in m};

    % Skalierung
    \foreach \x/\label in {0/0, 2/2, 4/4, 6/6, 8/8, 10/10} {
        \draw (\x,-0.1) -- (\x,0.1);
        \node[below] at (\x,-0.2) {\small \label};
    }
    \foreach \y/\label in {0/0, 1/100, 2/200, 3/300, 4/400, 5/500} {
        \draw (-0.1,\y) -- (0.1,\y);
        \node[left] at (-0.2,\y) {\small \label};
    }

    % Bewegung
    \draw[thick, physik] (0,0) -- (2,2);
    \draw[thick, physik] (2,2) -- (5,2);
    \draw[thick, physik] (5,2) -- (10,5);

    % Beschriftungen
    \node[loesung, above] at (1,1) {\scriptsize schnell};
    \node[loesung, above] at (3.5,2) {\scriptsize Stillstand};
    \node[loesung, above] at (7.5,3.5) {\scriptsize langsam};
\end{tikzpicture}
\end{center}

\begin{enumerate}[label=\alph*), leftmargin=2em]
    \item Beschreibung der Abschnitte: \hfill \lsg{3 P (je 1 P)}

    \textbf{Abschnitt 1} (0--2 min): \lsg{Lisa bewegt sich schnell / legt 200 m zurück / steile Gerade}

    \textbf{Abschnitt 2} (2--5 min): \lsg{Lisa steht still / Stillstand / waagerechte Linie / 0 m Weg}

    \textbf{Abschnitt 3} (5--10 min): \lsg{Lisa bewegt sich langsam / flache Gerade / legt 300 m zurück}

    \item Geschichte: \hfill \lsg{3 P}

    \textbf{Beispiellösungen:}

    \lsg{``Lisa fährt schnell mit dem Fahrrad los (200 m in 2 min).}
    \lsg{Dann muss sie an der Ampel warten (3 min Stillstand).}
    \lsg{Danach fährt sie langsamer weiter bis zur Schule (300 m in 5 min).''}

    \vspace{0.3em}
    \lsg{oder: ``Lisa joggt schnell los. Bei ihrem Freund wartet sie 3 Minuten.}
    \lsg{Dann gehen beide langsam gemeinsam zur Schule.''}

    \vspace{0.3em}
    \lsg{Bewertung: Drei Phasen erkannt (1 P) + sinnvolle Geschichte (1 P) + passt zum Diagramm (1 P)}
\end{enumerate}

\vspace{1em}

% ============================================================================
% PUNKTEÜBERSICHT
% ============================================================================
\hrule
\vspace{0.5em}

\begin{center}
\renewcommand{\arraystretch}{1.3}
\begin{tabular}{|l|c|c|c|}
\hline
\textbf{Niveau} & \textbf{Aufgaben} & \textbf{max. Punkte} & \textbf{Kompetenzen} \\
\hline
\stern{*} Basis & 1, 2, 3 & 20 P & Zuordnen, Berechnen, Ablesen \\
\hline
\stern{**} Standard & + 4, 5 & 33 P & + Zeichnen, Vergleichen \\
\hline
\stern{***} Erweitert & + 6 & 39 P & + Interpretieren, Transfer \\
\hline
\end{tabular}
\end{center}

\vspace{0.5em}

\noindent\textbf{Notenschlüssel:}

\begin{center}
\small
\begin{tabular}{|c|c|c|c|c|c|c|}
\hline
\textbf{Note} & 1 & 2 & 3 & 4 & 5 & 6 \\
\hline
\textbf{ab Punkte} & 37 (95\%) & 31 (80\%) & 23 (60\%) & 16 (40\%) & 8 (20\%) & 0 \\
\hline
\end{tabular}
\end{center}

\vspace{0.5em}

\noindent\textbf{AFB-Verteilung:}

\begin{center}
\small
\begin{tabular}{|l|c|c|l|}
\hline
\textbf{AFB} & \textbf{Punkte} & \textbf{Anteil} & \textbf{Aufgaben} \\
\hline
I (Reproduktion) & 14 P & 36\% & 1, 2a, 3 \\
\hline
II (Reorganisation) & 17 P & 44\% & 2b, 4, 5 \\
\hline
III (Transfer) & 8 P & 20\% & 6 \\
\hline
\end{tabular}
\end{center}

\vspace{0.5em}

\noindent\textbf{Rahmenplanbezug (MV Kl. 6 OS):}
\begin{itemize}[leftmargin=2em, topsep=2pt, itemsep=0pt]
    \item Einteilung nach Bahnform und Bewegungsart
    \item Weg und Zeit als physikalische Größen
    \item s(t)-Diagramme -- graphische Darstellung
\end{itemize}

\end{document}
